% GENERATED FILE. Edit canonical lessons, then run npm run generate:books.
% canonical-source-hash: fnv1a64:963682af61482b87
% canonical-lessons: TA-C300-avacaranilai, TA-C300-uratanku, TA-C300-mutanku, TA-C300-nilaikulai, TA-C300-itintuvilu

\chapter{A State of Emergency}
\label{ch:ta-a-state-of-emergency}
\hlchaptermodality{300}

\begin{chapteropening}
\textbf{By the end of this chapter:} \emph{I can say what happens to a town in an emergency.}

Complete the last lesson of chapter 300: I can say what happens to a town in an emergency.
\end{chapteropening}

\section[\texorpdfstring{\ta{அவசரநிலை} (avacaranilai)}{அவசரநிலை (avacaranilai)}]{\ta{அவசரநிலை} (avacaranilai) — an emergency}

\label{lesson:TA-C300-avacaranilai}



\begin{quote}
Before the new one: say the Tamil for a firefighter, then the Tamil for to suffocate.
\end{quote}

\subsection*{You'll want to know: \ta{அவசரநிலை}}
\textbf{\ta{அவசரநிலை}} — \emph{avacaranilai} — \textquotedblleft{}an emergency\textquotedblright{}.

\begin{quote}
\ta{மழை} \ta{நின்றால்}, \ta{அவசரநிலை} \ta{முடிவுக்கு} \ta{வரும்}.
\end{quote}

\begin{quote}
\emph{maḻai niṉṟāl, avacaranilai muṭivukku varum.}
\end{quote}

\begin{quote}
If the rain stops, the emergency will come to an end.
\end{quote}

From \ta{அவசரம்}, urgency, and \ta{நிலை}, a state.

\begin{cousinweb}
\textbf{\ta{அவசரநிலை}} (\emph{avacaranilai}) joins \ta{அவசரம்}, \textquotedblleft{}urgency\textquotedblright{}, borrowed from Sanskrit \emph{avasara}, \textquotedblleft{}occasion\textquotedblright{} (CDIAL 858), with native \ta{நிலை}, \textquotedblleft{}state, condition\textquotedblright{} (DEDR 3675).
\end{cousinweb}

\begin{grammarlens}[title={What you've built}]
One more word for telling what happened.
\end{grammarlens}

\subsection*{Your turn}
\begin{itemize}
\raggedright
  \item \emph{Say it:} \emph{avacaranilai}
  \item \emph{Say it:} \emph{avacaranilai}, once more
  \item \emph{Say it:} \emph{maḻai niṉṟāl, avacaranilai muṭivukku varum.} — If the rain stops, the emergency will come to an end.
  \item \emph{Recall:} say the Tamil for a firefighter, then the Tamil for to suffocate, then say \emph{avacaranilai} again
\end{itemize}

\begin{tcolorbox}[breakable,colback=teal!4,colframe=teal!35!black,title={Before you move on}]
What does \ta{அவசரநிலை} mean? (\textquotedblleft{}An emergency\textquotedblright{}.) Say it once more.
\end{tcolorbox}
\section[\texorpdfstring{\ta{ஊரடங்கு} (ūraṭaṅku)}{ஊரடங்கு (ūraṭaṅku)}]{\ta{ஊரடங்கு} (ūraṭaṅku) — a curfew}

\label{lesson:TA-C300-uratanku}



\begin{quote}
Before the new one: say the Tamil for to suffocate, then the Tamil for an emergency.
\end{quote}

\subsection*{You'll want to know: \ta{ஊரடங்கு}}
\textbf{\ta{ஊரடங்கு}} — \emph{ūraṭaṅku} — \textquotedblleft{}a curfew\textquotedblright{}.

\begin{quote}
\ta{நகரத்தில்} \ta{மூன்று} \ta{நாள்} \ta{ஊரடங்கு} \ta{இருந்தது}.
\end{quote}

\begin{quote}
\emph{nakarattil mūṉṟu nāḷ ūraṭaṅku iruntatu.}
\end{quote}

\begin{quote}
There was a three-day curfew in the city.
\end{quote}

\ta{ஊர்}, town + \ta{அடங்கு}, to settle down: the town goes quiet.

\begin{cousinweb}
\textbf{\ta{ஊரடங்கு}} (\emph{ūraṭaṅku}) is a Tamil compound built from \ta{ஊர்}, \textquotedblleft{}village, town\textquotedblright{} (DEDR 752), and \ta{அடங்கு}, \textquotedblleft{}to submit, settle down, cease\textquotedblright{} (DEDR 63).
\end{cousinweb}

\begin{grammarlens}[title={What you've built}]
One more word for telling what happened.
\end{grammarlens}

\subsection*{Your turn}
\begin{itemize}
\raggedright
  \item \emph{Say it:} \emph{ūraṭaṅku}
  \item \emph{Say it:} \emph{ūraṭaṅku}, once more
  \item \emph{Say it:} \emph{nakarattil mūṉṟu nāḷ ūraṭaṅku iruntatu.} — There was a three-day curfew in the city.
  \item \emph{Recall:} say the Tamil for to suffocate, then the Tamil for an emergency, then say \emph{ūraṭaṅku} again
\end{itemize}

\begin{tcolorbox}[breakable,colback=teal!4,colframe=teal!35!black,title={Before you move on}]
What does \ta{ஊரடங்கு} mean? (\textquotedblleft{}A curfew\textquotedblright{}.) Say it once more.
\end{tcolorbox}
\section[\texorpdfstring{\ta{முடங்கு} (muṭaṅku)}{முடங்கு (muṭaṅku)}]{\ta{முடங்கு} (muṭaṅku) — to stall; to grind to a halt}

\label{lesson:TA-C300-mutanku}



\begin{quote}
Before the new one: say the Tamil for an emergency, then the Tamil for a curfew.
\end{quote}

\subsection*{You'll want to know: \ta{முடங்கு}}
\textbf{\ta{முடங்கு}} — \emph{muṭaṅku} — \textquotedblleft{}to stall; to grind to a halt\textquotedblright{}.

\begin{quote}
\ta{வெள்ளத்தால்} \ta{நகரம்} \ta{முழுவதும்} \ta{முடங்கியது}.
\end{quote}

\begin{quote}
\emph{veḷḷattāl nakaram muḻuvatum muṭaṅkiyatu.}
\end{quote}

\begin{quote}
Because of the flood the whole city ground to a halt.
\end{quote}

Past (it …): \ta{முடங்கியது}, \emph{muṭaṅkiyatu}.

Said of a town, traffic or work that stops; the cause takes -\ta{ஆல்}.

\begin{cousinweb}
\textbf{\ta{முடங்கு}} (\emph{muṭaṅku}) is native Dravidian (DEDR 4919, 4933): to contract, to be held back, to stop moving.
\end{cousinweb}

\begin{grammarlens}[title={What you've built}]
One more word for telling what happened.
\end{grammarlens}

\subsection*{Your turn}
\begin{itemize}
\raggedright
  \item \emph{Say it:} \emph{muṭaṅku}
  \item \emph{Say it:} \emph{muṭaṅku}, once more
  \item \emph{Say it:} \emph{veḷḷattāl nakaram muḻuvatum muṭaṅkiyatu.} — Because of the flood the whole city ground to a halt.
  \item \emph{Recall:} say the Tamil for an emergency, then the Tamil for a curfew, then say \emph{muṭaṅku} again
\end{itemize}

\begin{tcolorbox}[breakable,colback=teal!4,colframe=teal!35!black,title={Before you move on}]
What does \ta{முடங்கு} mean? (\textquotedblleft{}To stall; to grind to a halt\textquotedblright{}.) Say it once more.
\end{tcolorbox}
\section[\texorpdfstring{\ta{நிலைகுலை} (nilaikulai)}{நிலைகுலை (nilaikulai)}]{\ta{நிலைகுலை} (nilaikulai) — to fall apart}

\label{lesson:TA-C300-nilaikulai}



\begin{quote}
Before the new one: say the Tamil for a curfew, then the Tamil for to stall.
\end{quote}

\subsection*{You'll want to know: \ta{நிலைகுலை}}
\textbf{\ta{நிலைகுலை}} — \emph{nilaikulai} — \textquotedblleft{}to fall apart\textquotedblright{}.

\begin{quote}
\ta{அந்தச்} \ta{செய்தியைக்} \ta{கேட்டு} \ta{நான்} \ta{நிலைகுலைந்தேன்}.
\end{quote}

\begin{quote}
\emph{antac ceytiyaik kēṭṭu nāṉ nilaikulaintēṉ.}
\end{quote}

\begin{quote}
I fell apart when I heard that news.
\end{quote}

Past (I …): \ta{நிலைகுலைந்தேன்}, \emph{nilaikulaintēṉ}.

From \ta{நிலை}, a steady state, and \ta{குலை}, to be upset.

\begin{cousinweb}
\textbf{\ta{நிலைகுலை}} (\emph{nilaikulai}) is a Tamil compound built from \ta{நிலை}, \textquotedblleft{}standing, a steady state\textquotedblright{} (DEDR 3675), and \ta{குலை}, \textquotedblleft{}to be upset, thrown into confusion\textquotedblright{} (DEDR 1806).
\end{cousinweb}

\begin{grammarlens}[title={What you've built}]
One more word for telling what happened.
\end{grammarlens}

\subsection*{Your turn}
\begin{itemize}
\raggedright
  \item \emph{Say it:} \emph{nilaikulai}
  \item \emph{Say it:} \emph{nilaikulai}, once more
  \item \emph{Say it:} \emph{antac ceytiyaik kēṭṭu nāṉ nilaikulaintēṉ.} — I fell apart when I heard that news.
  \item \emph{Recall:} say the Tamil for a curfew, then the Tamil for to stall, then say \emph{nilaikulai} again
\end{itemize}

\begin{tcolorbox}[breakable,colback=teal!4,colframe=teal!35!black,title={Before you move on}]
What does \ta{நிலைகுலை} mean? (\textquotedblleft{}To fall apart\textquotedblright{}.) Say it once more.
\end{tcolorbox}
\section[\texorpdfstring{\ta{இடிந்துவிழு} (iṭintuviḻu)}{இடிந்துவிழு (iṭintuviḻu)}]{\ta{இடிந்துவிழு} (iṭintuviḻu) — to collapse}

\label{lesson:TA-C300-itintuvilu}



\begin{quote}
Before the new one: say the Tamil for to stall, then the Tamil for to fall apart.
\end{quote}

\subsection*{You'll want to know: \ta{இடிந்துவிழு}}
\textbf{\ta{இடிந்துவிழு}} — \emph{iṭintuviḻu} — \textquotedblleft{}to collapse\textquotedblright{}.

\begin{quote}
\ta{பழைய} \ta{பாலம்} \ta{திடீரென்று} \ta{இடிந்துவிழுந்தது}.
\end{quote}

\begin{quote}
\emph{paḻaiya pālam tiṭīreṉṟu iṭintuviḻuntatu.}
\end{quote}

\begin{quote}
The old bridge suddenly collapsed.
\end{quote}

Past (it …): \ta{இடிந்துவிழுந்தது}, \emph{iṭintuviḻuntatu}.

Often written as two words: \ta{இடிந்து} \ta{விழு}.

\begin{cousinweb}
\textbf{\ta{இடிந்துவிழு}} (\emph{iṭintuviḻu}) is a Tamil compound built from \ta{இடிந்து}, \textquotedblleft{}having crumbled\textquotedblright{}, from \ta{இடி}, \textquotedblleft{}to crumble, fall to pieces\textquotedblright{} (DEDR 432), and \ta{விழு}, \textquotedblleft{}to fall\textquotedblright{} (DEDR 5430).
\end{cousinweb}

\begin{grammarlens}[title={What you've built}]
One more word for telling what happened.
\end{grammarlens}

\subsection*{Your turn}
\begin{itemize}
\raggedright
  \item \emph{Say it:} \emph{iṭintuviḻu}
  \item \emph{Say it:} \emph{iṭintuviḻu}, once more
  \item \emph{Say it:} \emph{paḻaiya pālam tiṭīreṉṟu iṭintuviḻuntatu.} — The old bridge suddenly collapsed.
  \item \emph{Recall:} say the Tamil for to stall, then the Tamil for to fall apart, then say \emph{iṭintuviḻu} again
\end{itemize}

\begin{tcolorbox}[breakable,colback=teal!4,colframe=teal!35!black,title={Before you move on}]
What does \ta{இடிந்துவிழு} mean? (\textquotedblleft{}To collapse\textquotedblright{}.) Say it once more.
\end{tcolorbox}
